% =====================================================================================
% Preface.
%
% NOT ONE NUMBER IN THIS FILE IS TYPED BY HAND: every figure in the prose is a macro
% from build/macros.tex, generated by cmp.report from the executed notebooks. The only
% literals are the decision thresholds (0.9 / 0.5) and the interval level (90 %), which
% are conventions stated in the text rather than measurements.
% =====================================================================================
\chapter*{Preface}
\addcontentsline{toc}{chapter}{Preface}
\markboth{Preface}{Preface}

This book is about one narrow and expensive question: \emph{did the marketing spend cause the
sales, and was it worth the money?} It is written for the analyst who has to answer it with the
data a firm actually holds --- a weekly sales panel, a campaign flag, a cost line, no coin flip
anywhere --- and who then has to defend the answer to someone holding the budget.

Two audiences are addressed at once, and the book was built so that neither has to read around
the other. The first is the MBA student at SDA Bocconi who will \emph{commission} this analysis
and must know what to demand of it: which question was actually answered, which assumption is
carrying the weight, and what the honest interval was. The second is the practising data
scientist who will \emph{run} it, and who wants the model, the diagnostics, the failure mode and
the code. The two need the same argument at different resolutions, so every chapter states its
decision in euros before it states its likelihood, and every chapter's technical machinery ends
by returning to the euros. What the book assumes is competence in ordinary regression --- a
reader who knows what a coefficient and a standard error are. What it assumes \emph{nothing}
about is causal inference: potential outcomes, confounding, backdoor paths and the do-operator
are built from scratch in \cref{ch:foundations}, and a reader who has met them before will lose
nothing by starting there anyway, because that chapter also plants the estimator ladder the rest
of the book climbs.

\section*{How to read it}

\Cref{ch:foundations,ch:pvp} should be read first, and in that order. They are the foundation:
the first supplies the vocabulary and the four estimators every later chapter falls back on, and
the second is the book's intellectual spine --- the exact difference between a confidence
interval and a posterior, and why only one of them can answer the question a budget-holder is
actually asking. Everything after them is self-contained. A reader who owns a geo experiment can
go straight to \cref{ch:geo}; a reader whose problem is a threshold-triggered loyalty tier can go
straight to \cref{ch:rdd}; the cross-references point back rather than forward, and the appendices
collect what the chapters share.

Each method chapter follows the same shape, and one feature of that shape is worth explaining
before it is met, because it is unusual and it is deliberate. \textbf{Every chapter estimates its
effect classically before it estimates it Bayesianly.} The classical baseline --- Step~0 --- comes
first: the simplest estimator that is still correct, a point estimate, and an interval whose
covariance assumption is stated out loud. No likelihood, no priors, no sampler. Only then does the
Bayesian model arrive, and it arrives with a burden of proof. The classical estimate is not a warm-up
exercise. It is the benchmark the posterior has to justify itself against, and the plain finding of
this book is that \emph{it frequently fails to}.

\section*{The argument}

That finding is worth stating in the preface rather than saving for a conclusion, because it is
the reason the book is organised the way it is. Four claims are made, and each is earned by
measurement in a chapter that could have refuted it.

\textbf{First: the two paradigms usually agree.} On the foundational example ---
\nbZeroNCustomers{} customers, one confounder, the same specification fitted twice --- ordinary
least squares and the posterior land \eur{\nbZeroEstGap} apart, and the 90\,\% credible interval
comes back \pc{\nbZeroWidthRatioPct} \emph{wider} than the confidence interval. Bayes did not
change the answer and bought nothing in precision. That is not an embarrassment; it is the
Bernstein--von Mises theorem behaving exactly as advertised, and it is the honest expectation to
carry into every chapter.

\textbf{Second: where the Bayesian arm lost, the likelihood was misspecified --- not the
paradigm.} \Cref{ch:geo} fits a synthetic control whose likelihood asserts that the weekly
errors are independent, applies it to a panel driven by a macroeconomic random walk, and
underprices the uncertainty of the post-launch total by a factor of \nbSevenUnderprice: its
nominal 90\,\%
interval covers the planted truth on \pc{\nbSevenTotalCov} of fresh panels, where the
design-based randomization interval covers on \pc{\nbSevenRiCov}. That result, alone, would read
as an indictment of Bayesian inference. It is not, and the proof is that \cref{ch:did} errs in
the \emph{opposite} direction --- its pooled posterior interval comes back
\nbEightBayesOverCluster$\times$ \emph{too wide}, because a large store baseline was left in the
residual variance. A paradigm cannot be systematically overconfident and systematically
underconfident at the same time. A \emph{specification} can be wrong in either direction, and a
specification is what differs between those two chapters. The diagnosis is therefore precise,
and so is the repair.

\textbf{Third: the prior does almost nothing; the likelihood is the whole ballgame.} Sweeping the
prior scale on the treatment effect across a factor of \nbZeroPriorRatio{} --- from a sceptic who
all but rules the effect out, to a prior flat across any plausible range --- moves the posterior
mean by \eur{\nbZeroPriorShift}. This is the answer to ``your priors are driving the result'',
and it has the merit of being a measurement rather than a reassurance. It is also the debugging
rule the book runs on: when a posterior disagrees with a design-based interval, do not adjust the
prior. Interrogate the likelihood --- and in particular interrogate what it assumed about the
correlation of the errors, the one assumption that is invisible in a good fit and fatal in a sum.

\textbf{Fourth: what survives everywhere is $\Prob(\text{effect} > \text{cost})$.} Every euro
decision in this book consumes that probability, and the classical arm structurally cannot produce
it --- not for want of effort, but because the framework that produced the confidence interval
contains no distribution over the effect whose mass above the cost line could be taken. The
posterior is in this book for that reason and, on the evidence assembled here, for very little
else. Which yields the discipline the book converges on, and the sentence it keeps returning to:
\textbf{use the posterior for the decision, and design-based inference for your confidence in it.}
Compute both. Where they agree, ship. Where they disagree, the likelihood is lying, and the
design-based number is the one to trust.

None of this is an argument against Bayesian modelling, and a reader who finishes the book
convinced that it was has read it backwards. It is an argument against a particular and very
common habit: reaching for a sampler in the belief that a posterior is a stronger claim than an
estimate. It is not a stronger claim. It is a \emph{different} claim, with a different price of
admission --- a model and a prior, both of which must then be defended --- and the only reliable
way to find out whether you have paid it is to have computed the classical answer first and to be
willing to be embarrassed by it.

\section*{What the book will not do}

It will not buy identification. This is the discipline every chapter enforces and
\cref{ch:pvp} demonstrates in its most uncomfortable form: drop one confounder from an otherwise
impeccable Bayesian model and the posterior concentrates, serenely and with excellent
diagnostics, on \eur{\nbZeroNoadjMean} against a truth of \eur{\nbZeroTrueDag} --- inside an
interval barely wider than the correct model's. Being confidently wrong is not punished with a
wide interval. Neither Bayesian inference nor machine learning has any opinion whatever about
whether the model it was handed identifies a causal effect. The graph decides what enters the
model; the likelihood and the prior only price the coefficients. Every chapter in this book
settles identification --- assumptions named, numbered, and either tested or explicitly declared
untestable --- \emph{before} it admires a posterior.
